% !TeX root = ../papers/nc-p3.tex
\section{階数15の分岐曲線の境界円板と有理第一座標}
\label{sec:rank-fifteen-boundary-disks}
\begin{theorem}\label{thm:rank-fifteen-boundary-disks}
二つの階数15分岐曲線上で、標数零の通常点を取る。101上のある付値で
$t$ の射影的還元が前節の非通常境界上にあるか、$\mathbb P^1(\mathbb F_{101})$
に属するとする。この点は標準Koszul左右AS正則で、各始点の道空間の直和は
四変数の四次中心多項式環上で階数256の自由加群である。
特に $t\in\mathbb Q$ の全通常点が含まれる。
\end{theorem}
\begin{proof}
まず分岐曲線上の四次中心の一般切断を使う。各二次関係では行ごとに、
中心元では四位置にわたり共通の分母を払う。通常点の長さ5の経路が定階数なので、
得られた切断は全通常点で中心である。無限遠では $x=1/t$、$y'=x^2y$ とし、
同じ単位で各行または一つの中心元の全四位置の極を払う。
この座標でも四次式は $F(x)=-x^4+14x^2-1$ である。

還元の境界を三通りに調べる。以下の根を含む二次体では共役な根を一度に扱う。
\begin{enumerate}
\item $t=0,1,-1,\infty$、$t^2+1=0$、$t^2\pm4t+1=0$ の各剰余円板で、
展開の中心を標数零で $\alpha$ とする。$F(\alpha)\ne0$ のとき
$u=t-\alpha$（無限遠では $u=x$）とし、$y$ の二根をHenselで持ち上げる。
線形因子では二根は共役である。$\alpha^2=-1$ では $y_0=\pm4\alpha$
を別々に取る。$F(\alpha)=0$ のときは $u=y$ として
$F(t)=u^2$ を $t=\alpha+\cdots$ について解く。全ての微分の分母は101で単元である。

定理\ref{thm:rank-fourteen-all-ordinary}の行消去をこの局所変数に適用する。
定数体は $\mathbb Q(\sqrt d)$、$d=-1,2,3$ のいずれかで、単元の検査には
二次ノルムを使う。単元ピボットがなければ残る全係数の定数項の両座標が
有理数として零であることを検査してから $u$ で割る。
通常な持上げでは $u\ne0$ である。精度は $t=\pm1$ で6、
$t=0$ で16、無限遠と $t^2+1=0$ で17、
二つの分岐二次因子で14および15から始める。
各段で四文字の最大の除算回数だけ精度を減らしても長さ13で正の精度が残る。

\item $t^2\pm2t-1=0$ の二因子では $\alpha=\pm1+\sqrt2$ と置き、
$u=t-\alpha$ とする。$\alpha=1+\sqrt2$ では $y_0=\pm2(\alpha+1)$、
$\alpha=-1+\sqrt2$ では $y_0=\pm2(1-\alpha)$ である。
$u$ の平方根を添加し、位置 $p$ の第 $j$ 矢を $u^{w_{pj}/2}$ 倍する。
各二次関係を、その語の係数の付値の最小値で規格化する。
例えば第一因子 $\alpha=1+\sqrt2$ の第一表示、$y_0=-2(\alpha+1)$ では
整数重みの四行を
\[
 (0,0,-1,-5),\quad(0,-4,-1,-1),\quad
 (0,0,-1,-5),\quad(0,4,3,3)
\]
と取れる。他の七組の重みも下記の検算に収めた。
$u^{12}$ を法とする二次体上の厳密な展開で各行の最低項を取り、
打切り後の項の重みが最小値より大きいことを検査する。
得られた最低項を $\mathbb F_{101^2}$ に還元すると、
各隣接関係の階数は6、経路次元は長さ4,5で35,56、四次中心は文字別に
$(3,2,0,0)$ である。この中心から取った四式で全四始点の長さ13まで
所定の商次元を得る。重み付き基底変換は付値体の有限拡大で可逆であり、
中心の持上げはこの新しい格子に適用できる。

\item 残る四次と八次の境界因子は101で六つの相異なる二次因子に分かれる。
それぞれの根と、その上の $y$ の両方に多項式切断を直接代入し、
所定の四式の商を長さ13まで検算する。Frobenius共役は同じ結果を与える。
\end{enumerate}
\path{computations/verify_rank_fifteen_boundary_disks.py} が以上を再現する。
第一と第三の場合では二つの連続する始点を調べ、
前節の二位置移動で残る二始点を得る。境界の和はこの移動で保たれる。
同一の中心四式は、各始点で有効な選択の非空開集合を交叉して取れる。

通常な有理剰余も全て調べられる。有限な非通常有理剰余は $0,1,-1,10,-10$ の五つである。
残り96個は $t\mapsto-1/t$ で48対となる。各対の代表において、
$F(t)$ が平方ならその両根、非平方なら二次体上の一根を試す。後者の他根はFrobenius共役である。
\path{computations/verify_rank_fifteen_rational_residues.py} は両表示・全四始点で
経路、中心および新しく選んだ四式の商を検算する。
いずれも長さ13で商が消える。一般点からのHilbert次元の下界と
命題\ref{prop:fixed-residual-regularity}の議論で全次数の結論を得る。
\end{proof}
